\section{Implementation}
\label{sec:impl}

The system is implemented in vLLM~v1 in roughly 3{,}900 lines across 22 files.
Every path is disabled by default and selected at runtime, so a build with the
feature off exercises the stock code; \S\ref{sec:eval:setup} reports what that
costs, because a baseline that is quietly slower than the engine it stands in
for would flatter every ratio in this paper.

\subsection{Host KV pool}
\label{sec:impl:pool}

The pool owns a private block-id space, a per-request map from logical block to
host block, and two pinned host tensors per layer. Host block ids are freed
only when a request finishes, which is what makes them safe: a request's device
blocks are released and reissued to other requests while it is still decoding,
so any mirror keyed by device block id would both lose the mapping needed to
find its own data and alias another request's (\S\ref{sec:design:spaces}).

Two layout decisions carry measurable weight. The mirrors are \textbf{block-major}
and $K$ and $V$ are \textbf{separate tensors}, so a run of consecutive host blocks is
one contiguous byte range on both sides and the per-layer transfer is a handful
of \texttt{cudaMemcpyAsync} calls issued straight from pinned memory. An earlier
layout that interleaved $K$ and $V$ made a run two chunks strided by a whole
plane; no pinned DMA was possible, every layer paid a host-side gather, and that
gather measured 22\,ms per layer --- 98\% of the streaming path and 72\% of
engine wall time. Restructuring the mirrors reduced it to 39\,ms in total.

That defect is worth reporting rather than quietly fixing, because it is
invisible to the metric one would naturally watch. It is a \emph{host} memcpy:
a faster interconnect does not remove it. On a coherent link the same code would
have spent hundreds of milliseconds per step gathering in order to hide a
transfer of a few milliseconds, and the natural conclusion would have been that
the interconnect does not help --- while measuring our own memory copies.

\subsection{Tier-aware attention}
\label{sec:impl:attn}

Attention is the only operation that distinguishes the tiers
(\S\ref{sec:design:exec}). We subclass the FlashAttention backend's forward
path. Device-tier rows are served by the standard call; host-tier rows are
gathered with \texttt{index\_select}, attended against the rolling slot with the
compacted block table, and scattered back into the shared output with
\texttt{index\_copy\_}.

The standard call still runs over every row, so host-tier rows would be attended
twice --- once uselessly, over their full history --- if nothing prevented it.
We set their \texttt{seqused\_k} to 1, so they read a single key in the first
call and are overwritten by the second. On A100 that single change was worth
22\% of wall time.

\subsection{Scheduler}
\label{sec:impl:sched}

Three changes make a host-tier request hold no device blocks
(\S\ref{sec:design:additive}): admission with an empty block list, a prompt
marked already computed so no prefill is scheduled, and deferral of a
block-starved request at the head of the waiting queue so it does not block
admissions that need no blocks at all. The third was found only by testing; the
first two are insufficient without it.

Reclaiming a resident request's blocks on demotion required repairing the
promotion path, and the failure mode was not a crash. Block accounting is
length-based, so a nulled entry still counts as present: \texttt{allocate\_slots}
would never refill it, and the promotion gate --- ``do the blocks cover the whole
sequence?'' --- would have \emph{passed on an all-null table} and served
garbage. Refill is therefore all-or-nothing, the gate counts non-null entries,
and refilled ids reach the worker through a rewrite channel, because they
replace positions rather than appending.

\subsection{Prefill handover}
\label{sec:impl:stash}

A host-tier request's KV must reach the host pool before it is admitted. We
implement a prompt-keyed stash: the producer writes KV into the pool, the entry
outlives the producing request, and a later request preloads from it. Admission
is contingent on the worker acknowledging that the KV is present, since the
scheduler cannot observe a host-side write and admitting on the assumption
would leave a request with no KV anywhere.

The handover costs \textbf{0.22\,ms per decode step} at 2048 context
(14.2\,ms amortised over 64 steps, 38\,GB/s), about 5\% of a host-tier
request's per-step cost. Two corrections produced that number and both are
worth recording. The path was first written block by block, which at 2048
context is 8{,}192 small copies per request; since the pool allocates
per-request contiguous extents, a layer is one copy, and fixing it gained
$2.2\times$. The first measurement was then taken on a \emph{contended} node
and reported 4.06\,GB/s; exclusive on an idle node the same code reports
38\,GB/s. \textbf{A $9.4\times$ error from contention alone}, in the measurement
that answers the most obvious objection to our system model --- which is why
\S\ref{sec:eval:setup} requires exclusive access for every absolute we report.

\subsection{What the streaming path costs, and why}
\label{sec:impl:cost}

Beyond the transfer itself, using the host path at all costs a fixed amount per
step. We reduced it from 37.0 to 12.3\,ms, and the residual is not a defect.

Four sites accounted for the reduction and all were the same species:
loop-invariant work executed inside the per-layer loop. A device reduction and
a \texttt{.item()} evaluated once per layer inside a guard whose purpose was to
\emph{skip} work --- an optimisation paying a full pipeline drain to decide
whether to optimise. Scatter indices rebuilt per layer from data fixed once the
step's schedule is known. A sequence maximum and an \texttt{arange} likewise.
A clone and index-fill of the masked sequence lengths. Each is invariant once
the step is scheduled; together they were worth 21\,ms per step, and the
end-to-end throughput deficit moved $-11.3\% \to -5.8\% \to -2.8\% \to -2.3\%$
as they were removed, with the control reproducing to within 0.1\,tokens/s
across five runs.

\paragraph{Method, because it decided the outcome twice.} Aggregate per-phase
timers attributed the cost wrongly, twice, and in both cases plausibly. A timer
wrapped around asynchronous GPU work charges waiting-for-the-GPU to whichever
phase holds it; a profiler that synchronises \emph{inside} its own timing
wrapper makes double-buffering impossible by construction and then reports that
double-buffering does not help. What located all four sites was an
operator-level trace differencing a step with one host-tier request against a
step with none. We record this because the same instinct --- read the phase
timings, form a hypothesis --- is what a reader would try first.

\paragraph{The residual is a floor, not a bug.} An operator trace against the
current code shows \texttt{cudaStreamSynchronize} gone from the table entirely
and no dominant remaining term: the cost is roughly 640 additional kernel
launches per step, about ten operations across 64 layers, at 0.19\,ms per
layer. That is the per-layer structure itself. Removing it requires batching
across layers rather than issuing per layer --- the rolling slots already exist
to support it --- which is a design change we have not made and do not claim.
It is also machine-independent, so it does not distort the comparison between
interconnects in \S\ref{sec:eval:cost}.
